  \documentclass[a4paper,12pt]{article}
  \usepackage{styles}

  \begin{document}

  \litSubsubsection{3.1.2 Optimal Device Occupancy Calculation}

  \subsubsection*{Hardware Constraints}

  Achievable occupancy is bounded by hardware limits that vary by compute capability. NVIDIA defines occupancy as ``the ratio of the number of active warps per multiprocessor to the maximum number of possible active warps''~\cite{nvidia2024cuda}. For Compute Capability 8.6, the CUDA Programming Guide specifies a maximum of 1536 threads per SM, which corresponds to 48 warps of 32 threads each~\cite{nvidia2024programming}. Each SM can host up to 16 concurrent thread blocks, and each block may contain at most 1024 threads. Together, these parameters define the limits within which kernel launch configurations must operate.

  \textbf{Occupancy Calculation.} Theoretical occupancy for a given kernel configuration is computed as:

  \begin{equation}
      \text{Occupancy} = \frac{\text{Active Threads per SM}}{\text{Maximum Threads per SM}} \times 100\%
  \end{equation}

  The number of active threads depends on the launch configuration (threads per block and blocks per grid) and on per-thread resource usage such as registers and shared memory. As noted in Section 2.2.2, this value is an upper bound; achieved occupancy may be lower due to workload imbalance or insufficient parallelism~\cite{nvidia2024occupancy}.

  \subsubsection*{Thread Block Configuration}

  The thread block size for fitness evaluation kernels must balance occupancy with memory access efficiency. The implementation uses a heuristic based on population size, divided into three regimes. For small populations ($N < 256$), the thread count per block is rounded up to the nearest multiple of the warp size (32), since the CUDA Programming Guide notes that ``thread blocks are best specified to have a total number of threads which is a multiple of 32,'' as partial warps leave hardware lanes unused~\cite{nvidia2025cudaguide}. Medium populations ($256 \leq N < 10{,}000$) use 256 threads per block, a value within the 128--256 range recommended by NVIDIA for general-purpose kernels~\cite{nvidia_occupancy}. Large populations ($N \geq 10{,}000$) increase the block size to 512 threads to reduce scheduling overhead while staying within the 1024-thread hardware limit.

  The number of blocks per grid is calculated as:

  \begin{equation}
      \text{Blocks per Grid} = \left\lceil \frac{N}{\text{Threads per Block}} \right\rceil
  \end{equation}

  This ensures the grid covers the entire population with one thread per individual. For the RTX 2050 with 16 SMs, the minimum population for full SM coverage at 256 threads per block is $16 \times 256 = 4{,}096$ individuals. The Section~3 benchmarks use population sizes between 50 and 1{,}000, so the block size is fixed at 256 threads throughout, consistent with the medium-population regime above.

  \subsubsection*{Population Size and Occupancy}

  Total GPU thread capacity influences the optimal population size. The implementation targets 50--100\% occupancy and also considers GA-specific heuristics that suggest population sizes of 10--50 times the problem dimensionality~\cite{katoch2021review}. Population sizes are aligned to multiples of 256 so that all warps in each block are fully populated.

  \textbf{Practical Constraints.} High occupancy is generally desirable but does not guarantee the best performance. As the CUDA Best Practices Guide states, ``higher occupancy does not always equate to higher performance''~\cite{nvidia2024cuda}. Volkov shows that memory-bound kernels can plateau at moderate occupancy levels, while instruction-level parallelism may compensate for lower occupancy in compute-bound cases~\cite{volkov2010occupancy}. GA fitness evaluation involves repeated memory accesses to population arrays and problem-specific data (distance matrices, training patterns), so the trade-off between occupancy and per-thread resources must be tested empirically. Sections 3.3 and 3.2 benchmarks explore these configurations across a range of population sizes.

  \end{document}